% tex/sections/conclusion.tex
% ──────────────────────────────────────────────────────────────────────────────
\chapter{Conclusion}
\label{chap:conclusion}
% ──────────────────────────────────────────────────────────────────────────────
%
% TODO: Write conclusion.
%
% Suggested structure:
%   - Restate the objectives
%   - Summarise key findings (no new content)
%   - Broader significance / take-home message

% Opening: restate research objective in one sentence
% (~2 sentences)

The aim of this thesis was to contribute to the understanding of seismic signatures of crack propagation and avalanche breakoff. This was done
by setting up a moving-source model of a crack propagating at sub-rayleigh and supershear speeds in snow using spectral element modelling in SALVUS. The results of the simple two-layer model were evaluated to determine how the seismic signature measured by a DAS cable is influenced by the crack propagation regime. To make the physics more realisitic, the cohesive zone model for mixed-mode crack propagation was set up. This was then followed by a model which translated the results of material point method simulations by \cite{trottetTransitionSubRayleighAnticrack2022} into seismic moment tensors. The results of this simulation were evaluated to verify effects observed in the simple model at more realistic conditions. 

\section*{Answer to research questions}

% RQ1: FE modelling framework
% One paragraph: SALVUS spectral element method, two-layer and
% five-layer geometry, custom STF from MPM, DAS via pyber gauge length.
% One sentence on key assumptions (elastic, homogeneous, 2D, prescribed speed).

% RQ2: Physical phenomena in simple model
% One paragraph per main finding:
%   - Two-wavefront structure in sub-Rayleigh (S0 precursor + crack front)
%   - Single merged front in supershear (vr ≈ vp)
%   - A0 Lamb mode dominance (polarity reversal evidence)
%   - Arrest signatures (fan vs V-shape)
%   - Amplitude difference (10× more energy in supershear)

% RQ3: MPM-guided results
% One paragraph:
%   - Sub-Rayleigh confirmed at realistic parameters (vr ≈ 57 m/s)
%   - Quasi-static strain field (DC component, absent from simple model)
%   - Five-layer geometry validates simple model signatures

% RQ4: Implications for DAS detection
% One paragraph (your most important result):
%   - Sub-Rayleigh may be undetectable at snow bottom
%   - Arrest signature distinguishes regimes
%   - False conclusion risk if only supershear phase detected
%   - Quasi-static field as potential early warning

\section*{Limitations}

% Bullet or short paragraph covering:
%   - 2D geometry (no along-slope propagation, no azimuth variation)
%   - Elastic only (no attenuation, no viscous effects)
%   - Prescribed rupture speed (not self-consistent)
%   - Homogeneous layers (no spatial variability in weak layer)
%   - No field validation of simulated DAS signatures

\section*{Recommendations for future work}

% Point to outlook chapter or summarise in 3-4 sentences:
%   - 3D extension for arbitrary cable orientation
%   - Self-consistent rupture speed from cohesive zone model
%   - Field data comparison (PST with DAS)
%   - Heterogeneous weak layer properties
